\chapter{Аналитический раздел}

\section{Постановка задачи}

Разработать загружаемый модуль для ОС Linux, позволяющий использовать стилус графического планшета в качестве мыши через Type-C соединение. Движение стилуса по графическому планшету должны преобразовываться в относительные перемещения курсора мыши, касания стилусом планшета - в нажатие левой кнопки мыши, а нажатия на кнопку стилуса - в нажатия правой кнопки мыши.

\section{Представление ввода стилуса}

В абсолютной модели значение координаты обозначает положение инструмента в рабочей области. В относительной модели сообщение содержит изменение положения~\cite{kernelinput}. Для перехода между ними требуется хранить предыдущий пригодный отсчёт и вычислять разность по каждой оси.

Первый отсчёт после появления пера сохраняется без формирования движения: прежнее положение может относиться к другой последовательности. При выходе из зоны считывания начало отсчёта сбрасывается. Это позволяет переносить перо между участками поверхности без скачка курсора при возвращении.

Для настройки чувствительности используется целочисленный делитель. Если дробная часть отбрасывается независимо на каждом шаге, малые перемещения теряются. Поэтому остатки деления накапливаются отдельно по осям и учитываются при следующем отсчёте. При сбросе начала последовательности остатки также обнуляются.

Контакт наконечника и нажатие боковой кнопки задают состояния. Пока наконечник касается поверхности, левая кнопка остаётся нажатой; движение в этом состоянии обеспечивает перетаскивание. При наведении движение разрешено без нажатия. Основная боковая кнопка аналогично управляет правой кнопкой мыши.

Состояния левой и правой кнопок независимы: боковая кнопка может быть нажата одновременно с контактом наконечника, и преобразователь должен сохранять эту комбинацию.

Выход пера из зоны считывания должен освобождать обе кнопки. Такое сообщение обрабатывается до проверки координат, поскольку координатные поля при выходе могут быть непригодны. Иначе отбрасывание пакета могло бы сохранить нажатую кнопку. Другие повреждённые пакеты игнорируются без частичного изменения предыдущего отсчёта.

\section{Особенности обмена по USB}

Физическое устройство может содержать несколько интерфейсов, например интерфейс пера и интерфейс совместимости с мышью. USB-драйвер привязывается к интерфейсу, поэтому выбор определяется его назначением и идентификаторами устройства~\cite{kernelusb}. Захват одного интерфейса не отключает остальные автоматически.

Конечная точка входящего направления выбирается по дескрипторам активной конфигурации. Её характеристики определяют тип передачи и допустимый размер буфера. Сравнение типов передачи приведено в таблице~\ref{tbl:usb-transfer}~\cite{usb_spec}.

\begin{table}[htbp]
    \caption{Выбор типа USB-передачи}
    \label{tbl:usb-transfer}
    \centering
    \small
    \begin{tabular}{|p{0.20\linewidth}|p{0.34\linewidth}|p{0.35\linewidth}|}
        \hline
        Тип передачи & Назначение & Применение в работе \\ \hline
        Управляющая & Настройка устройства & Установка режима протокола \\ \hline
        Прерываний & Периодически обслуживаемые сообщения & Приём состояния пера \\ \hline
        Потоковая & Передача массивов данных & Не соответствует выбранной точке ввода \\ \hline
        Изохронная & Потоки с временными ограничениями & Не требуется для событий стилуса \\ \hline
    \end{tabular}
\end{table}

Обслуживание входящей конечной точки организует USB-хост. Название «передача прерываний» не означает произвольного инициирования обмена периферией. Интервал обслуживания характеризует транспорт, но не определяет полную задержку отображения результата пользователю~\cite{usb_spec}.

Асинхронная операция в Linux описывается USB Request Block (URB). Запрос содержит ссылку на устройство, буфер, параметры конечной точки и обработчик завершения. После выполнения доступны статус и фактическое число принятых байтов. Для продолжения приёма URB отправляется повторно, пока работа разрешена~\cite{kernelurb}.

Размер USB-буфера нельзя считать длиной сообщения. До обращения к координатам проверяются фактическая длина и идентификатор отчёта. Входной отчёт содержит текущие данные, а HID-дескриптор описывает их формат. Использование известной раскладки пакета в разборщике не требует подмены дескриптора~\cite{usb_hid}.

Для интерфейса с загрузочным протоколом и протоколом отчётов необходимо выбрать режим, соответствующий разборщику. Стандартный управляющий запрос устанавливает режим обмена; при ошибке запуск обработки с неподтверждённым форматом прекращается.

\section{Сравнение способов преобразования}

Пользовательская программа может получать события существующего устройства и создавать мышь через uinput. Подход использует готовый разбор аппаратных данных и позволяет отлаживать преобразование вне ядра~\cite{uinput}. Однако он требует постоянно работающего процесса, управления доступом к устройству и подавления исходных событий. Алгоритм при этом выполняется вне разрабатываемого модуля ядра.

Специализированный HID-драйвер использует штатный USB-транспорт. Разделение транспорта и интерпретации отчётов сокращает объём собственного кода, но оставляет доставку пакетов компоненту \texttt{usbhid}~\cite{hidtransport}. Данный подход допустим для многих задач настройки ввода, однако не реализует выбранную в работе непосредственную обработку USB-запросов.

Собственный USB-драйвер занимает согласованные интерфейсы, принимает байты через URB и формирует стандартные события мыши~\cite{kernelusb}. Он требует самостоятельного управления повторными запросами, ошибками и временем жизни ресурсов. Это усложнение соответствует цели системного программирования: реализовать собственную компоненту обработки внешнего устройства на уровне ядра.

Итог сравнения приведён в таблице~\ref{tbl:approaches}. Соответствие задаче оценивается с учётом размещения преобразования в ядре и непосредственного управления USB-обменом.

\begin{table}[htbp]
    \caption{Сравнение способов обработки ввода}
    \label{tbl:approaches}
    \centering
    \small
    \begin{tabular}{|p{0.25\linewidth}|p{0.29\linewidth}|p{0.35\linewidth}|}
        \hline
        Способ & Собственная обработка & Ограничение для работы \\ \hline
        Процесс и uinput & Преобразование готовых событий & Выполнение вне ядра \\ \hline
        HID-драйвер & Интерпретация отчётов & Штатный USB-транспорт \\ \hline
        USB-драйвер & Приём и преобразование пакетов & Полное управление URB и ресурсами \\ \hline
        Подмена дескриптора & Изменение описания полей & Запрещена требованиями \\ \hline
    \end{tabular}
\end{table}

Для дальнейшей разработки выбран собственный USB-драйвер. Аппаратная специализация ограничивается выбором интерфейсов и разбором согласованного формата. Алгоритм преобразования и передача результата подсистеме ввода отделяются от этих особенностей.

Выбор уровня USB обосновывается составом разрабатываемой компоненты, а не предполагаемым преимуществом быстродействия. Сравнение задержек пользовательского процесса и драйвера требует отдельного измерения на одинаковом оборудовании. Без такого эксперимента меньший объём промежуточной обработки нельзя считать доказательством меньшей задержки управления.

\section{Передача событий и надёжность}

Результат должен быть независим от исходного аппаратного пакета. Драйвер регистрирует устройство Linux Input Subsystem с относительными осями \texttt{REL\_X}, \texttt{REL\_Y} и кнопками \texttt{BTN\_LEFT}, \texttt{BTN\_RIGHT}. Событие синхронизации объединяет изменения одного отчёта в пакет~\cite{kernelinput}. Абсолютные координаты и давление наружу не передаются.

Далее данные получают штатные потребители ввода. В графическом сеансе обработка может включать libinput и компоненты оконной системы. Прямое управление экранным указателем из драйвера не требуется~\cite{libinput}. Такое разделение сохраняет независимость от конкретного приложения.

Если интерфейс совместимости продолжает выдавать исходные события, одновременно действуют два способа управления указателем. Поэтому соответствующий интерфейс временно занимается без создания второго устройства ввода. Уже подключённый штатный драйвер не вытесняется одной загрузкой модуля: необходима явная перепривязка выбранных интерфейсов. Предыдущие привязки сохраняются для восстановления после выгрузки; остальные устройства не затрагиваются.

Асинхронный обмен требует согласовать обработчик завершения и остановку. Сначала запрещаются новые отправки, затем ожидается завершение URB и отложенного восстановления, после чего освобождаются ресурсы~\cite{ldd3,love3}. Проверка разрешения работы и повторная отправка защищаются совместно, чтобы отключение не могло вклиниться между ними.

Обработчик завершения выполняет короткую обработку пакета. Восстановление конечной точки, требующее ожидания, переносится в рабочую очередь~\cite{kernelurb}. Штатная отмена запроса не должна запускать новый обмен. Если восстановить приём невозможно, кнопки освобождаются, а причина записывается в журнал. Отдельный контекст каждого устройства исключает смешение координат при нескольких подключениях.

Ошибка формата и ошибка транспорта имеют разные последствия. Непригодное содержимое успешно принятого сообщения не требует перенастройки USB-точки: пакет отбрасывается, приём продолжается. Ошибка обмена обрабатывается по статусу завершения. Такое разделение сохраняет предыдущий корректный отсчёт при нарушении формата и не смешивает его с восстановлением аппаратного канала.

\clearpage
